\documentclass{article}
\usepackage{amsmath,amssymb,amsthm}
\usepackage[noend]{algorithm2e}
\usepackage{hyperref}
\usepackage{enumitem}
\newtheorem{theorem}{Theorem}
\newtheorem{lemma}{Lemma}
\newtheorem{assumption}{Assumption}
\newcommand{\EE}{\mathbb{E}}
\newcommand{\RR}{\mathbb{R}}
\newcommand{\norm}[1]{\left\lVert #1\right\rVert}
\newcommand{\inner}[2]{\left\langle #1,#2\right\rangle}
\begin{document}

\section*{Algorithm Name}
Differentially Private Stochastic Gradient Descent with Gaussian Noise (DP‑SGD)

\section*{Mathematical Formulation}
Let $f:\RR^{d}\rightarrow\RR$ be the (possibly non‑convex) loss function.  
At iteration $t=0,1,\dots,T-1$ we perform

\[
\boxed{
\begin{aligned}
g_{t} &:= \nabla f(w_{t})+\xi_{t}, \qquad 
\xi_{t}\sim\mathcal N(0,\sigma^{2}I_{d}) ,\\[4pt]
w_{t+1} &:= w_{t}-\eta_{t}\,g_{t},
\end{aligned}}
\]

where 
\begin{itemize}[noitemsep]
\item $w_{t}\in\RR^{d}$ is the model parameter,
\item $\eta_{t}>0$ is the learning‑rate (may be constant or decreasing),
\item $\sigma>0$ controls the magnitude of the added Gaussian noise.
\end{itemize}
The sequence $\{w_{t}\}$ is released only after the final iteration; the Gaussian mechanism together with a standard moments accountant guarantees $(\varepsilon,\delta)$‑differential privacy for a given sampling scheme.

\section*{Pseudocode}
\begin{algorithm}[H]
\caption{DP‑SGD with Gaussian Noise}
\KwIn{Initial point $w_{0}\in\RR^{d}$, learning‑rate schedule $\{\eta_{t}\}_{t=0}^{T-1}$, noise scale $\sigma>0$, number of iterations $T$}
\For{$t=0$ \KwTo $T-1$}{
    Compute the (full‑batch) gradient $ \nabla f(w_{t}) $\;
    Sample $\xi_{t}\sim\mathcal N(0,\sigma^{2}I_{d})$\;
    Set $g_{t}\leftarrow \nabla f(w_{t})+\xi_{t}$\;
    Update $w_{t+1}\leftarrow w_{t}-\eta_{t}g_{t}$\;
}
\Return $w_{T}$
\end{algorithm}

\section*{Dry Run Example}
Consider the one‑dimensional quadratic $f(w)=(w-3)^{2}$, start at $w_{0}=0$, use a constant step size $\eta=0.1$, noise variance $\sigma^{2}=0$ (so we can illustrate the deterministic part) and run three iterations.

\begin{center}
\begin{tabular}{c|c|c|c}
$t$ & $\nabla f(w_{t})$ & $g_{t}= \nabla f(w_{t})+\xi_{t}$ & $w_{t+1}=w_{t}-\eta g_{t}$ \\ \hline
0 & $2(w_{0}-3)=-6$ & $-6$ & $0-0.1(-6)=0.6$ \\[4pt]
1 & $2(w_{1}-3)=2(0.6-3)=-4.8$ & $-4.8$ & $0.6-0.1(-4.8)=1.08$ \\[4pt]
2 & $2(w_{2}-3)=2(1.08-3)=-3.84$ & $-3.84$ & $1.08-0.1(-3.84)=1.464$ 
\end{tabular}
\end{center}
Corresponding objective values:
$f(w_{0})=9,\;f(w_{1})=(0.6-3)^{2}=5.76,\;f(w_{2})=(1.08-3)^{2}=3.6864,\;f(w_{3})=(1.464-3)^{2}=2.3689.$
The example shows the classic geometric decrease for a smooth convex quadratic; the same update rule with $\xi_{t}\neq0$ yields a private trajectory.

\newpage
\section*{Assumptions}
\begin{assumption}[Smoothness]\label{ass:smooth}
$f$ is $L$‑smooth, i.e. for all $x,y\in\RR^{d}$
\[
f(y)\le f(x)+\inner{\nabla f(x)}{y-x}+\frac{L}{2}\norm{y-x}^{2}.
\]
\end{assumption}
\begin{assumption}[Bounded Gradient Variance]\label{ass:var}
For every $t$, the noise $\xi_{t}$ satisfies
\[
\EE[\xi_{t}]=0,\qquad \EE[\norm{\xi_{t}}^{2}]=d\sigma^{2},
\]
and it is independent of $w_{t}$ and of all past randomness.
\end{assumption}
\begin{assumption}[Learning‑rate]\label{ass:lr}
The step size is constant $\eta_{t}\equiv\eta$ with 
\[
0<\eta\le\frac{1}{L}.
\]
\end{assumption}

\section*{Main Convergence Theorem}
\begin{theorem}\label{thm:main}
Assume \ref{ass:smooth}–\ref{ass:lr}. Let $\{w_{t}\}_{t=0}^{T}$ be generated by the DP‑SGD update. Then
\[
\frac{1}{T}\sum_{t=0}^{T-1}\EE\!\left[\norm{\nabla f(w_{t})}^{2}\right]
\;\le\;
\frac{2\bigl(f(w_{0})-f^{\*}\bigr)}{\eta T}
\;+\;
L\eta d\sigma^{2},
\]
where $f^{\*}=\inf_{w}f(w)$. Consequently, choosing $\eta = \Theta\!\bigl(T^{-1/2}\bigr)$ yields
\[
\min_{0\le t<T}\EE\!\left[\norm{\nabla f(w_{t})}^{2}\right]=O\!\left(\frac{1}{\sqrt{T}}+ \sigma^{2}\right).
\]
\end{theorem}

\section*{Theoretical Properties}
\begin{itemize}
\item \textbf{Stability.} The bound contains the term $L\eta d\sigma^{2}$, showing that larger noise (higher privacy) inflates the stationary gradient norm linearly in $\sigma^{2}$.
\item \textbf{Hyper‑parameter sensitivity.} The optimal constant step size balances the $1/(\eta T)$ term against $L\eta d\sigma^{2}$; the minimizer is $\eta^{\star}= \sqrt{2\bigl(f(w_{0})-f^{\*}\bigr)/(L d\sigma^{2}T)}$.
\item \textbf{Comparison to non‑private SGD.} Setting $\sigma=0$ recovers the classic $O(1/(\eta T))$ bound for smooth non‑convex SGD.
\item \textbf{Privacy‑utility trade‑off.} The additive term $L\eta d\sigma^{2}$ is precisely the utility loss due to the Gaussian mechanism; decreasing $\sigma$ (less privacy) improves the bound.
\end{itemize}

\section*{Discussion}
The theorem shows that DP‑SGD converges to a point where the expected gradient norm is $O(T^{-1/2})$ up to an additive noise floor proportional to $\sigma^{2}$. This matches the best known rates for private non‑convex optimization and demonstrates that the algorithm is stable under the smoothness assumption. The result also guides the choice of $\eta$ and $\sigma$ for a prescribed privacy budget.

\newpage
\section*{Preliminaries (Definitions and Lemmas)}
\begin{lemma}[Smoothness inequality]\label{lem:smooth}
If $f$ is $L$‑smooth then for any $x,y\in\RR^{d}$,
\[
f(y)\le f(x)+\inner{\nabla f(x)}{y-x}+\frac{L}{2}\norm{y-x}^{2}.
\]
\end{lemma}
\begin{proof}
Standard consequence of the Lipschitz continuity of $\nabla f$. \hfill $\square$
\end{proof}
\begin{lemma}[Unbiasedness of the noisy gradient]\label{lem:unbiased}
For the DP‑SGD gradient $g_{t}= \nabla f(w_{t})+\xi_{t}$,
\[
\EE[g_{t}\mid w_{t}]=\nabla f(w_{t}),\qquad
\EE\!\left[\norm{g_{t}}^{2}\mid w_{t}\right]=\norm{\nabla f(w_{t})}^{2}+d\sigma^{2}.
\]
\end{lemma}
\begin{proof}
$\EE[\xi_{t}]=0$ and $\EE[\norm{\xi_{t}}^{2}]=d\sigma^{2}$ by Assumption~\ref{ass:var}. Independence of $\xi_{t}$ from $w_{t}$ yields the identities. \hfill $\square$
\end{proof}

\section*{Main Proof (Step‑by‑Step Derivation)}
We prove Theorem~\ref{thm:main}. All expectations are taken over the randomness of $\{\xi_{t}\}_{t=0}^{T-1}$ unless otherwise noted.

\bigskip
\noindent\textbf{[Step 1]} Apply Lemma~\ref{lem:smooth} with $x=w_{t}$ and $y=w_{t+1}=w_{t}-\eta g_{t}$:
\[
f(w_{t+1})\le f(w_{t})+\inner{\nabla f(w_{t})}{-\,\eta g_{t}}+\frac{L}{2}\norm{-\,\eta g_{t}}^{2}.
\tag{1}
\]

\noindent\textbf{[Step 2]} Rearrange (1):
\[
f(w_{t+1})\le f(w_{t})-\eta\inner{\nabla f(w_{t})}{g_{t}}+\frac{L\eta^{2}}{2}\norm{g_{t}}^{2}.
\tag{2}
\]

\noindent\textbf{[Step 3]} Expand the inner product using $g_{t}= \nabla f(w_{t})+\xi_{t}$:
\[
\inner{\nabla f(w_{t})}{g_{t}}=\norm{\nabla f(w_{t})}^{2}+\inner{\nabla f(w_{t})}{\xi_{t}}.
\tag{3}
\]

\noindent\textbf{[Step 4]} Substitute (3) into (2):
\[
f(w_{t+1})\le f(w_{t})-\eta\norm{\nabla f(w_{t})}^{2}
-\eta\inner{\nabla f(w_{t})}{\xi_{t}}
+\frac{L\eta^{2}}{2}\norm{g_{t}}^{2}.
\tag{4}
\]

\noindent\textbf{[Step 5]} Take conditional expectation $\EE[\cdot\mid w_{t}]$ on both sides of (4). Using $\EE[\inner{\nabla f(w_{t})}{\xi_{t}}\mid w_{t}=0$ (zero‑mean noise) we obtain
\[
\EE\!\left[f(w_{t+1})\mid w_{t}\right]\le f(w_{t})-\eta\norm{\nabla f(w_{t})}^{2}
+\frac{L\eta^{2}}{2}\EE\!\left[\norm{g_{t}}^{2}\mid w_{t}\right].
\tag{5}
\]

\noindent\textbf{[Step 6]} Apply Lemma~\ref{lem:unbiased} to the last term:
\[
\EE\!\left[\norm{g_{t}}^{2}\mid w_{t}\right]
= \norm{\nabla f(w_{t})}^{2}+d\sigma^{2}.
\tag{6}
\]

\noindent\textbf{[Step 7]} Insert (6) into (5):
\[
\EE\!\left[f(w_{t+1})\mid w_{t}\right]\le
f(w_{t})-\eta\norm{\nabla f(w_{t})}^{2}
+\frac{L\eta^{2}}{2}\Bigl(\norm{\nabla f(w_{t})}^{2}+d\sigma^{2}\Bigr).
\tag{7}
\]

\noindent\textbf{[Step 8]} Collect the terms that involve $\norm{\nabla f(w_{t})}^{2}$:
\[
\EE\!\left[f(w_{t+1})\mid w_{t}\right]\le
f(w_{t})-\Bigl(\eta-\frac{L\eta^{2}}{2}\Bigr)\norm{\nabla f(w_{t})}^{2}
+\frac{L\eta^{2}}{2}d\sigma^{2}.
\tag{8}
\]

\noindent\textbf{[Step 9]} By Assumption~\ref{ass:lr}, $\eta\le 1/L$, hence $\eta-\frac{L\eta^{2}}{2}\ge \frac{\eta}{2}$. Therefore
\[
\EE\!\left[f(w_{t+1})\mid w_{t}\right]\le
f(w_{t})-\frac{\eta}{2}\norm{\nabla f(w_{t})}^{2}
+\frac{L\eta^{2}}{2}d\sigma^{2}.
\tag{9}
\]

\noindent\textbf{[Step 10]} Take total expectation on (9) and rearrange:
\[
\frac{\eta}{2}\EE\!\left[\norm{\nabla f(w_{t})}^{2}\right]
\le \EE\!\left[f(w_{t})\right]-\EE\!\left[f(w_{t+1})\right]
+\frac{L\eta^{2}}{2}d\sigma^{2}.
\tag{10}
\]

\noindent\textbf{[Step 11]} Sum (10) over $t=0,\dots,T-1$:
\[
\frac{\eta}{2}\sum_{t=0}^{T-1}\EE\!\left[\norm{\nabla f(w_{t})}^{2}\right]
\le f(w_{0})-\EE\!\left[f(w_{T})\right]
+ \frac{L\eta^{2}}{2}Td\sigma^{2}.
\tag{11}
\]

\noindent\textbf{[Step 12]} Since $f(w_{T})\ge f^{\*}$, replace $\EE[f(w_{T})]$ by $f^{\*}$ to obtain an upper bound:
\[
\frac{\eta}{2}\sum_{t=0}^{T-1}\EE\!\left[\norm{\nabla f(w_{t})}^{2}\right]
\le f(w_{0})-f^{\*}
+ \frac{L\eta^{2}}{2}Td\sigma^{2}.
\tag{12}
\]

\noindent\textbf{[Step 13]} Divide both sides of (12) by $\eta T$:
\[
\frac{1}{T}\sum_{t=0}^{T-1}\EE\!\left[\norm{\nabla f(w_{t})}^{2}\right]
\le
\frac{2\bigl(f(w_{0})-f^{\*}\bigr)}{\eta T}
+ L\eta d\sigma^{2}.
\tag{13}
\]
Equation (13) is exactly the bound claimed in Theorem~\ref{thm:main}. \hfill $\square$

\section*{Rate Derivation (With Constants)}
From (13) we see two competing terms:
\[
A(\eta)=\frac{2\Delta}{\eta T},\qquad
B(\eta)=L\eta d\sigma^{2},\quad\text{where }\Delta:=f(w_{0})-f^{\*}.
\]
Minimising $A(\eta)+B(\eta)$ w.r.t. $\eta>0$ yields
\[
\eta^{\star}= \sqrt{\frac{2\Delta}{L d\sigma^{2}T}}.
\]
Plugging $\eta^{\star}$ back gives
\[
\frac{1}{T}\sum_{t=0}^{T-1}\EE\!\left[\norm{\nabla f(w_{t})}^{2}\right]
\le 2\sqrt{\frac{2L d\sigma^{2}\Delta}{T}}.
\]
Thus the average squared gradient norm decays as $O\!\bigl(T^{-1/2}\bigr)$ plus the unavoidable noise floor $L\eta d\sigma^{2}=O(\sigma\sqrt{d/T})$.

\section*{Special Cases}
\begin{enumerate}
\item \textbf{No privacy (}\boldmath$\sigma=0$\textbf{).} The bound reduces to
\[
\frac{1}{T}\sum_{t=0}^{T-1}\EE\!\left[\norm{\nabla f(w_{t})}^{2}\right]\le \frac{2\Delta}{\eta T},
\]
the classic non‑convex SGD rate.
\item \textbf{Small learning rate.} If $\eta\le 1/(L)$ is taken much smaller than the optimal $\eta^{\star}$, the $L\eta d\sigma^{2}$ term dominates and the algorithm stagnates at a gradient norm of order $L\eta d\sigma^{2}$.
\item \textbf{Large batch (deterministic gradient).} The analysis still holds because the only stochasticity stems from the Gaussian noise; the proof does not rely on sampling variance.
\end{enumerate}

\end{document}